\section{仿真结果}
\subsection{与适配 GT-ML-PDA 的补充比较}
在同一批阈后观测上，\gtmethod 的平均 GOSPA 比 \method 高 56.00 m（表~\ref{tab:gt}）。该结论在二维、已知五成员刚性槽位条件下建立。

消融结果尚不能区分幅度标记与关联边缘化各自的独立贡献，本文将性能归因于整体机制。
